\section{Failure Modes, Security, and Operational Control}
\label{sec:failure}

The preceding sections examined what KG--LLM systems can do. This section examines how they break, and which engineering controls contain the damage. The organizing claim is that a deployed KG--LLM system is a versioned configuration whose knowledge substrate mutates, which makes configuration management, lineage, and adversarial robustness first-order technical concerns rather than deployment paperwork.

\subsection{Semantic mapping as a model-risk source}

Interoperability is often reduced to ``the system uses FHIR''~\cite{hl7fhir} or ``it maps to OMOP''~\cite{omop}. Neither claim tells you whether meaning survived the mapping, and meaning is what the graph is for. Mappings from local terms to SNOMED CT, ICD, RxNorm, LOINC, UMLS, or OMOP concepts can be one-to-many, approximate, version-dependent, or context-sensitive. Every graph node derived from a patient record should therefore retain the source code and text, the mapped concept, the mapping method and confidence, the terminology version, and the effective time. Unmapped and ambiguous concepts must stay visible, because silently dropping them produces false reassurance, and unit conversion, negation, family history, uncertainty, and temporal intervals all need explicit handling.

The testable consequence is that conformance is not correctness. Interoperability testing should measure clinical meaning through cross-site phenotype agreement, concept-coverage and ambiguity rates, unit and temporal consistency, preservation of source qualifiers, round-trip tests, and performance stratified by mapped versus unmapped concepts. A schema-valid message can still be clinically wrong, and a conformant analytic table can still reflect a lossy extraction pipeline.

\subsection{The deployed system is a versioned configuration}

The intervention is not ``the LLM.'' It is the configuration
\begin{equation}
C = (M,\ P,\ G,\ T,\ R,\ V,\ I),
\end{equation}
where $M$ is the model and version, $P$ the prompt and tool policy, $G$ the KG snapshot and schema, $T$ the terminology and ontology releases, $R$ the retrieval and reranking policy, $V$ the verification and abstention layer, and $I$ the interface and workflow integration. A change to any element produces a new configuration unless equivalence is demonstrated.

This formulation closes a reproducibility gap that version-controlling the model alone leaves open. Updating a drug database, swapping an embedding model, adding an ontology synonym, altering a hop limit, or changing where a result appears in the interface can all change behavior while $M$ stays fixed. Release documentation should therefore carry immutable component identifiers, graph build manifests, source checksums, evaluation results, and rollback targets, and every logged output should name the configuration that produced it. Without this, a regression cannot be attributed and a rollback cannot be verified.

\subsection{Provenance, contradiction, and time as representation problems}

Biomedical facts are rarely timeless binary relations. An edge may come from a randomized trial, an observational association, a preclinical experiment, a database inference, a text-extraction model, or an LLM proposal, and its validity may be population-specific, dose-specific, disputed, superseded, or retracted. A graph intended to ground clinical claims should represent at least source identity, publication or record date, extraction method, evidence type, population and experimental qualifiers, confidence, review status, and valid time. MDKG shows what this buys, attaching conditional statements, baseline population attributes, source information, and confidence to relations~\cite{gao2025mdkg}.

Two representation choices follow. First, contradictions should be preserved rather than collapsed into a single consensus edge, so that retrieval can return supporting, opposing, and insufficient evidence and generation can report the disagreement together with the date through which evidence was searched. Second, retraction has to propagate. Removing one document from an index is not enough, because dependent edges, community summaries, cached embeddings, generated rationales, and downstream derived assertions all need to be identified and invalidated or rebuilt. This is a lineage problem, and solving it is what lets an operator answer, after an incident, which outputs used the affected evidence and which configuration should replace them.

\subsection{Hazard analysis across the lifecycle}

Failure-mode analysis becomes actionable only when each hazard links to a control and to objective verification evidence. Table~\ref{tab:hazards} instantiates this for the lifecycle of Fig.~\ref{fig:lifecycle}. Severity is task-dependent: a missing literature citation and a missed drug contraindication do not warrant the same control budget, so the table should be re-instantiated per intended use rather than applied uniformly.

\begin{table*}[!t]
\centering
\caption{\textbf{How a graph-grounded system fails at each stage, and what evidence would show the control is working.} Listing hazards is easy and cheap; the columns that matter are the last two. A control without verification evidence is an intention, so each row pairs a stage of the lifecycle in Fig.~\ref{fig:lifecycle} with something measurable. The stages are ordered by when they occur, and the failures propagate downward: a bad edge admitted at ingestion becomes a retrieval candidate, a retrieval gap becomes an omitted contraindication, and an unreviewed write-back becomes tomorrow's apparently independent evidence. Severity is task-dependent and the table should be re-instantiated per intended use, because a missing literature citation and a missed drug interaction do not warrant the same control budget.}
\label{tab:hazards}
\footnotesize
\begin{tabularx}{\textwidth}{@{}L{3.0cm} L{2.6cm} Y Y@{}}
\toprule
\textbf{Stage and hazard} & \textbf{Consequence} & \textbf{Required controls} & \textbf{Verification evidence} \\
\midrule
Source ingestion: outdated, retracted, or malicious source & False or obsolete assertions & Source allowlist and quality tiers; retraction feed; checksums; independent-source rule; quarantine & Source audit; retraction drill; \% edges with complete provenance; rebuild test \\
Construction: entity or relation error & Wrong traversal and downstream output & Schema validation; terminology constraints; calibrated extraction; dual-model or human review for high-risk edges & Relation precision/recall by type; calibration; inter-rater agreement; sampled edge audit \\
Patient mapping: negation, experiencer, date, unit error & Wrong patient state and unsafe personalization & Assertion and temporal model; unit validation; source-text trace; ambiguity retention & Mapping coverage; temporality/negation test set; cross-site audit; expert adjudication \\
Retrieval: incomplete or irrelevant subgraph & Omitted contraindication or misleading context & Hybrid retrieval; minimum sufficiency; conflict retrieval; context budget; abstention & Recall@$k$; path validity; contraindication-recall stress test; oracle comparison \\
Generation: unsupported inference & Hallucinated diagnosis or recommendation & Claim decomposition; graph and text verification; constrained output; severity-aware abstention & Unsupported-claim rate; citation entailment; counterfactual and distractor tests \\
Human use: automation bias or alert fatigue & Inappropriate acceptance or workarounds & Role-based actions; calibrated uncertainty; escalation paths & Override and verification rates; time, load, and missed-error analysis \\
Closed-loop update: proposal re-enters corpus & Circular self-validation and error amplification & Candidate-edge quarantine; independent corroboration; immutable commits; no self-citation & Lineage query; simulated circularity test; reversible update drill \\
Operations: drift, outage, unapproved change & Silent degradation or uncontrolled behavior & Monitoring; configuration registry; safe fallback; rollback & Drift dashboard; failover exercise; configuration attestation; rollback recovery time \\
\bottomrule
\end{tabularx}
\end{table*}

\subsection{Circular self-validation}

Bidirectional systems carry a failure mode that unidirectional ones do not. An LLM proposes an edge, the edge is committed to the graph, and a later LLM retrieves it as though it were independent evidence. Repetition across agents or model calls does not create independent confirmation; the system has cited its own derivative output, and the citation can be syntactically correct while being epistemically circular. Self-consistency checks and agreement among closely related models amplify the effect rather than detecting it, because correlated errors look like corroboration.

The control is a lifecycle for candidate knowledge that is distinct from curated knowledge. Proposed edges enter a quarantine zone carrying explicit generated-by, source, model, prompt, and confidence metadata. Promotion requires evidence not derived from the same model output, a policy-defined source threshold, schema checks, and human approval for high-impact relation types. Downstream retrieval excludes candidate edges by default or labels them unambiguously, graph commits are immutable and reversible, and an edge derived from an LLM-generated summary retains lineage to the original article, because the summary is not a second source. Agent tools should hold least privilege: retrieval agents get read access, extraction agents may write candidates, and only an authorized release process writes asserted edges.

\subsection{Privacy as a graph-structural problem}

Patient-specific graphs can be more disclosive than the records they derive from. Linking rare diagnoses, dates, locations, family relations, genomic variants, and social factors eases re-identification even after direct identifiers are removed, and the attack surface extends past the graph itself to embeddings, vector indexes, cached prompts, retrieved subgraphs, logs, and human feedback, all of which belong in the data inventory. Legal obligations vary by jurisdiction and deployment, with HIPAA~\cite{hhsocr} and the GDPR~\cite{gdpr} setting the applicable duties in the United States and the European Union. The engineering point is that neither is satisfied by a compliance label attached to a system whose graph leaks structure.

The architectural response follows from that observation. Keep the global biomedical graph separate from patient graphs, bind access by purpose and role to the smallest subgraph that answers the question, and authorize at query time rather than trusting a downstream filter. Everything derived from a patient graph inherits its sensitivity, so embeddings, cached context, and logs need the same retention limits and egress controls as the records themselves. The full architectural checklist is in the supplementary material~\cite{supplementary}. Federated deployment reduces central data movement without guaranteeing privacy, because model updates, gradients, graph statistics, and query patterns can each leak, so a privacy-enhancing technology should be evaluated against a stated threat model rather than adopted as a label.

\subsection{Security and adversarial robustness}

Graph grounding changes the attack surface rather than removing it. Yang et al.~\cite{yang2024poisoning} showed that a single maliciously generated abstract could substantially shift drug--disease rankings in a literature-derived medical KG. Alber et al.~\cite{alber2025poisoning} showed that poisoning a small fraction of medical training content raised harmful output while leaving standard benchmark performance largely unchanged. PoisonedRAG showed that a few injected documents can steer RAG answers~\cite{zou2025poisonedrag}. Together these results refute the assumption that retrieval from an external corpus is intrinsically safer than parametric memory, and they explain why benchmark accuracy is not a safety measurement.

A KG--LLM threat model spans four planes. The \emph{knowledge plane} covers poisoned publications, compromised databases, false triples, ontology manipulation, entity hijacking, malicious summaries, and stale evidence. The \emph{control plane} covers direct and indirect prompt injection, malicious instructions embedded in retrieved text, unsafe query translation, tool-policy bypass, and unauthorized graph updates. The \emph{data plane} covers cross-patient retrieval, inference from embeddings, log leakage, overbroad traversal, and insecure backups. The \emph{action plane} covers unvalidated output passed to orders, messages, scheduling, code execution, or other agents, along with excessive permissions and resource-exhaustion attacks.

Controls layer accordingly. Ingestion uses source allowlists, cryptographic manifests, provenance thresholds, anomaly detection, and quarantine. Retrieved content is treated as untrusted data and kept separate from system instructions, so that a document cannot redefine tool permissions. Generated graph queries are parsed against an allowlisted grammar, executed read-only with least privilege, bounded by time and result size, and checked for protected-data scope. Agent tools carry explicit capability boundaries, confirmation for consequential actions, egress restrictions, and complete audit logs. High-risk claims are verified against an independently curated source set rather than the graph that generated them. Security evaluation should report attack success rate and residual harm, not accuracy after perturbation, and safe failure means abstention or a bounded evidence response rather than a silent fallback to ungrounded model knowledge. The NIST Generative AI Profile offers a cross-sector structure for organizing this work across the lifecycle~\cite{nistgenai2024}.

\subsection{Monitoring a system whose knowledge mutates}

A single accuracy measure is too slow and too coarse to catch KG--LLM degradation, because the causes sit at different stages and a drop at one stage can be masked by slack at another. Monitoring has to be stage-local. The signals that matter divide along the same lines as the evaluation layers: whether inputs are still mapping to concepts, whether the graph is still fresh and internally consistent, whether retrieval is still returning evidence, whether generation is still supported by what came back, and whether users are still behaving as though they trust it appropriately. A worked signal list is in the supplementary material~\cite{supplementary}.

What matters more than the list is what happens when a signal moves. Thresholds should be set in advance and should trigger something specific: investigation, rollback, suspension, or a restricted mode. Sentinel cases should be rerun after every material change to $C$, so that a series of individually harmless updates cannot drift the system without anyone noticing. For a dynamic evidence graph, a high-risk system should retrieve from approved immutable snapshots and promote a new snapshot only after regression testing, because retrieving from a live graph means the system's behavior changes without a release.

Incident response needs traceability in both directions. From a harmful output, an investigator has to reach the patient inputs, mappings, retrieved edges, model, policy, and user actions that produced it. From a corrupted source, they have to reach every release and every output that depended on it. Systems that can do the first but not the second can explain a single failure and cannot bound its blast radius.

\subsection{Deployment obligations beyond engineering}

Engineering controls do not exhaust what a deployment requires. A KG--LLM system used in care falls under medical-device regulation when its intended purpose makes it a device, and under data-protection, professional, and institutional duties regardless. The relevant instruments include the WHO guidance on large multimodal models in health~\cite{who2024lmm}, FDA guidance on clinical decision support software and on lifecycle management for AI-enabled device software together with the predetermined change control framework~\cite{fdacds2026,fdalifecycle2025,fdapccp2025}, the IMDRF and FDA good machine learning practice principles~\cite{gmlp2025}, and the EU AI Act~\cite{euaiact2024}. Reporting guidance is similarly established, spanning TRIPOD-LLM and TRIPOD+AI for development and evaluation, PROBAST+AI for risk of bias, STARD-AI for diagnostic accuracy, DECIDE-AI for early live evaluation, SPIRIT-AI and CONSORT-AI for trials, and FUTURE-AI for deployment~\cite{gallifant2025tripodllm,collins2024tripodai,moons2025probastai,sounderajah2025stardai,vasey2022decideai,rivera2020spiritai,liu2020consortai,lekadir2025futureai}.

Two implications matter for an engineering reader. The configuration $C$ is what these frameworks govern, so a change to the graph snapshot or the retrieval policy can constitute a regulated change even when the model weights are untouched. And responsibility is distributed across source licensors, graph curators, model providers, system integrators, health systems, and clinical users, so no single actor can certify the whole chain. Detailed accountability allocation and jurisdiction-specific classification lie outside the scope of this survey; the clinical-informatics reviews compared in Section~\ref{sec:priorsurveys} treat them directly~\cite{murali2026integrating,liu2025biomedrag}.
